% Generated draft: requires mathematical and source review.
\documentclass[11pt]{article}
\usepackage[margin=25mm]{geometry}
\usepackage{fontspec}
\setmainfont{Latin Modern Roman}
\usepackage{amsmath,amssymb,mathtools,unicode-math}
\setmathfont{Latin Modern Math}
\usepackage{enumitem,needspace,xurl,hyperref}
\setlength{\parindent}{0pt}
\setlength{\parskip}{5pt}
\setlength{\emergencystretch}{3em}
\newcommand{\N}{\mathbb{N}}
\newcommand{\Z}{\mathbb{Z}}
\newcommand{\Q}{\mathbb{Q}}
\newcommand{\R}{\mathbb{R}}
\newcommand{\C}{\mathbb{C}}
\newcommand{\F}{\mathbb{F}}
\newcommand{\A}{\mathbb{A}}
\newcommand{\PP}{\mathbb{P}}
\newcommand{\E}{\mathbb{E}}
\newcommand{\eps}{\varepsilon}
\newcommand{\dd}{\,\mathrm{d}}
\begin{document}
\section*{UnsolvedMath ID: 20000418}
\textbf{Source problem code:} catalogue entry
\subsection{Definitions and mathematical statement}
\textbf{Original Problem 7 (K. Murty).} A Hilbert modular surface in characteristic $p$, attached to a real quadratic field $F$, has cycles parametrizing Hilbert--Blumenthal abelian varieties with prescribed Newton polygons. Are their classes in the span of special cycles consisting of Hecke translates of Hirzebruch--Zagier curves, and how can one check this? The source does not specify the prime behavior, level, compactification or equivalence relation on cycle classes.

The original comments ask: are the cycles pulled back from the Siegel threefold $A_2$ (``not always'', with A. Langer's work pertinent, J. Achter and C. Schoen); could McMullen's Teichmüller curves generate exotic classes (J. Ellenberg); can van der Geer and Moonen's material make the characteristic-$p$ cycles explicit (J. Achter); can one intersect these cycles with reductions of the CM Tate cycles in Murty's lecture (J. Ellenberg); and how does the nontrivial $\theta\in\operatorname{Gal}(F/\Q)$ act (J. Getz)? The lecture's cycles and these references are not supplied in detail.

\textbf{Attached unreviewed inert-prime analysis.} Restrict to an inert prime $p>2$ and hyperspecial good reduction. The aid identifies the two labelled Newton divisor families as Goren--Oort partial-Hasse divisors $D_1,D_2$. Their sum is reported to be the geometric pullback of the Siegel Hasse divisor; each labelled divisor has nonzero $\theta$-anti-invariant Néron--Severi component and is not individually such a pullback. The level and chosen surface must admit the stated involution for that eigenspace formulation. The summary does not give all compactification or boundary conventions.

In a cuspidal $\pi$-component with distinct Satake parameters, the aid attributes to Tian--Xiao's matrix two independent special-fiber Tate directions. It compares these with characteristic-zero potential-Tate dimension $t$: reductions of HZ/twisted-HZ cycles and their prime-to-$p$ Hecke translates are reported to miss both directions if $t=0$, and at least one if $t\le1$. The exact automorphic hypotheses and matrix are not supplied. This is about reductions of the stated characteristic-zero cycles, not every possible special-fiber algebraic cycle.

The aid's finite containment criterion is numerical. On a smooth projective model, choose a $\theta$-invariant ample class $H$ and use the primitive Néron--Severi space $H^\perp$ with its negative-definite intersection pairing. After removing the known tautological part of the Goren--Oort component span, the aid says the residual quotient is generated by component-augmentation images (differences of components in each labelled family modulo the totals). For a $\theta$-stable proposed special-cycle span, work separately in the even and odd primitive blocks. Let $S_1,\ldots,S_r$ be an independent basis of the relevant proposed span, $D$ a residual class in that block, and
\[A=(S_i\cdot S_j)_{i,j},\qquad b=(D\cdot S_i)_i.\]
The reported exact test is
\[D\in\operatorname{span}_{\Q}\{S_1,\ldots,S_r\}\quad\Longleftrightarrow\quad D^2-b^{\mathsf T}A^{-1}b=0.\]
The matrix must be nonsingular, which follows for an independent basis in the stated negative-definite primitive space. Numerical class coefficients, intersection numbers, a complete proposed basis and the tautological projection are required inputs, not supplied by the criterion. This does not decide Chow-theoretic containment or construct the suggested CM/Teichmüller intersections. All analysis remains confined to the aid's stated inert good-reduction setting and requires review.
\subsection{Short English statement}
Murty asks whether Newton-polygon cycles on Hilbert modular surfaces come from Hirzebruch--Zagier special cycles, with five related geometric questions. The aid treats inert good primes, reports new special-fiber Tate directions beyond certain reductions, and gives a conditional intersection-matrix test for numerical containment in each involution eigenspace.
\subsection{Sources}
\begin{itemize}
\item Local UnsolvedMath registry, permanent ID 20000418, source code catalogue entry; \texttt{aim-algebraic-number-theory-notes.json}, notes[69], tag \texttt{problem}.
\item AIM workshop \emph{The Tate conjecture}, July 23--27, 2007, Problem 7 (K. Murty), list transcribed by Christopher Lyons: \url{https://aimath.org/WWN/tateconjecture/tateconjecture.pdf}. The supplied verification identifies ``2contained'' as a page-number artifact; no PDF was consulted here.
\item Original comments and attributions: J. Achter, C. Schoen, J. Ellenberg and J. Getz; references by name to A. Langer, McMullen, van der Geer and Moonen, and Murty's lecture. No publication or lecture details are supplied.
\item Attached machine-generated research attempt 1 refers to Goren--Oort divisors and Tian--Xiao's matrix without publication locators. Result \texttt{partial\_result}, status \texttt{partially\_solved}, rigor label \texttt{proved}, novelty confidence low. It is not peer reviewed and requires independent verification before citation.
\end{itemize}
\end{document}
